% Options for packages loaded elsewhere
\PassOptionsToPackage{unicode}{hyperref}
\PassOptionsToPackage{hyphens}{url}
\documentclass[
  ukrainian,
]{article}
\usepackage{xcolor}
\usepackage[margin=2.5cm]{geometry}
\usepackage{amsmath,amssymb}
\setcounter{secnumdepth}{-\maxdimen} % remove section numbering
\usepackage{iftex}
\ifPDFTeX
  \usepackage[T1]{fontenc}
  \usepackage[utf8]{inputenc}
  \usepackage{textcomp} % provide euro and other symbols
\else % if luatex or xetex
  \usepackage{unicode-math} % this also loads fontspec
  \defaultfontfeatures{Scale=MatchLowercase}
  \defaultfontfeatures[\rmfamily]{Ligatures=TeX,Scale=1}
\fi
\usepackage{lmodern}
\ifPDFTeX\else
  % xetex/luatex font selection
\fi
% Use upquote if available, for straight quotes in verbatim environments
\IfFileExists{upquote.sty}{\usepackage{upquote}}{}
\IfFileExists{microtype.sty}{% use microtype if available
  \usepackage[]{microtype}
  \UseMicrotypeSet[protrusion]{basicmath} % disable protrusion for tt fonts
}{}
\makeatletter
\@ifundefined{KOMAClassName}{% if non-KOMA class
  \IfFileExists{parskip.sty}{%
    \usepackage{parskip}
  }{% else
    \setlength{\parindent}{0pt}
    \setlength{\parskip}{6pt plus 2pt minus 1pt}}
}{% if KOMA class
  \KOMAoptions{parskip=half}}
\makeatother
\usepackage{graphicx}
\graphicspath{{/Users/kalinicenkovitalijmikolajovic/.codex/worktrees/5c15/Script_for_translate_book/Output_Final/farlou_chunked_v1_pdf_chunks/}}
\makeatletter
\newsavebox\pandoc@box
\newcommand*\pandocbounded[1]{% scales image to fit in text height/width
  \sbox\pandoc@box{#1}%
  \Gscale@div\@tempa{\textheight}{\dimexpr\ht\pandoc@box+\dp\pandoc@box\relax}%
  \Gscale@div\@tempb{\linewidth}{\wd\pandoc@box}%
  \ifdim\@tempb\p@<\@tempa\p@\let\@tempa\@tempb\fi% select the smaller of both
  \ifdim\@tempa\p@<\p@\scalebox{\@tempa}{\usebox\pandoc@box}%
  \else\usebox{\pandoc@box}%
  \fi%
}
% Set default figure placement to htbp
\def\fps@figure{htbp}
\makeatother
\ifLuaTeX
\usepackage[bidi=basic,shorthands=off,provide=*]{babel}
\else
\usepackage[bidi=default,shorthands=off,provide=*]{babel}
\fi
\ifLuaTeX
  \usepackage{selnolig} % disable illegal ligatures
\fi
\setlength{\emergencystretch}{3em} % prevent overfull lines
\providecommand{\tightlist}{%
  \setlength{\itemsep}{0pt}\setlength{\parskip}{0pt}}
\usepackage{bookmark}
\IfFileExists{xurl.sty}{\usepackage{xurl}}{} % add URL line breaks if available
\urlstyle{same}
\hypersetup{
  pdflang={uk},
  hidelinks,
  pdfcreator={LaTeX via pandoc}}

\author{}
\date{}

\begin{document}

\section{ФОРМУЛА
ДАЛАМБЕРА}\label{ux444ux43eux440ux43cux443ux43bux430-ux434ux430ux43bux430ux43cux431ux435ux440ux430}

ЦЕЛЬ ЛЕКЦИИ: Построить решение задачи Коши для волнового уравнения

и
\emph{\{yy\}=\texttt{u\_\{l+1,\textbackslash{}\ j\}}-2u}\{l,~j\}+\texttt{u\_\{l-1,\textbackslash{}\ j\}},
(центральные разностные производные),

и подставим их в уравнение с частными производными. Соответствующая
сетка показана на рис. 43.2. Разрешив получившееся уравнение
относительно \texttt{u\_\{i,l\}} , получаем

Посмотрим внимательно на это уравнение. Коэффициенты при
\texttt{u\_\{i+1,\textbackslash{}\ j\},\textbackslash{}\ u\_i,\textbackslash{}\ j+i,\textbackslash{}\ u\_l,\textbackslash{}\ j-1}
н \texttt{u\_\{i-1\},\textbackslash{}\ j} положительны, и их сумма равна
единице. Другими словами, решение \texttt{u\_i,\textbackslash{}\ j}
является взвешенным средним решений в четырех соседних точках.
Следовательно, можно Р И С

\pandocbounded{\includegraphics[keepaspectratio]{images/_page_324_Picture_9.jpeg}}

. 43.2. Узлы сетки для случайного маршруга.

модифицировать «Блуждающего пьяницу» так, чтобы вероятности перехода в
соседние точки были не 0,25, а равнялись бы коэффициенту в
соответствующем члене. Другими словами, если пьяница находится в точке
(i, j), то он переходит в точку

Эта задача описывает движение бесконечной струны с заданными начальными
условиями. Она была решена в 1750 году французским математиком
Даламбером. Формула Даламбера

с вероятностью
\texttt{\textbackslash{}frac\{1\}\{2(1+\textbackslash{}sin\ x\_j)\}} ,
\texttt{(i,\ j-1)} с вероятностью
\texttt{\textbackslash{}frac\{1\}\{2(1+\textbackslash{}sin\ x\_j)\}} ,
\texttt{(i+1,\ j)} с вероятностью
\texttt{\textbackslash{}frac\{\textbackslash{}sin\ x\_j\}\{2(1+\textbackslash{}sin\ x\_j)\}}
, \texttt{(i-1,\ j)} с вероятностью
\texttt{\textbackslash{}frac\{\textbackslash{}sin\ x\_j\}\{2(1+\textbackslash{}sin\ x\_j)\}}
.

В остальном игра не меняется. Модификация для других задач может
оказаться более хитроумной, но основная идея остается той же. Читатель
может заняться разработкой своих игр для решения других задач.
Параболический случай рассмотрен в разделе «Задачи».

\section{ЗАМЕЧАНИЯ}\label{ux437ux430ux43cux435ux447ux430ux43dux438ux44f}

\texttt{1.\ Отметим,\ что\ если\ относительные\ частоты\ попадания\ \ \textquotesingle{}P\_A(p\_l)\textquotesingle{}\ \ известны,\ то\ можно\ легко\ найти\ величину\ \ \textquotesingle{}u\textbackslash{}left(A\textbackslash{}right)\textquotesingle{}\ \ при\ других\ граничных\ условиях\ \ \textquotesingle{}g\_l\textquotesingle{}\ .\ Для\ этого\ достаточно\ подставить\ \ \textquotesingle{}g\_l\textquotesingle{}\ \ в\ формулу}

позволяет легко найти решение, если известны начальные условия. Кроме
того, она позволяет дать интересную физическую интерпретацию решения на
языке бегущих волн.

\begin{verbatim}
Как, наверное, помнит читатель, в параболническом случае мы сначала решали диффузионные задачи на конечном отрезке (методом разделения переменных), а уже затем перешли к неограниченным промежуткам (—   \infty < x < \infty  ), где применяли преобразование Фурье. В гиперболическом случае мы поступим наоборот. Начнем с решения одномерного волнового уравнения на всей действительной оси, т. е. с задачи Коши
\end{verbatim}

(17.1)

а не разыгрывать вновь случайные прогулки.

\begin{enumerate}
\def\labelenumi{\arabic{enumi}.}
\setcounter{enumi}{1}
\tightlist
\item
  Во многих случаях исследователя интересует решение уравнения с
  частными производными в одной точке. Если граница области сложна и
  уравнение зависит от трех или четырех независимых переменных, метод
  Монте-Карло может оказаться удобным. Методы Монте-Карло возникли при
  решении очень трудных задач диффузии нейтронов. Эти задачи невозможно
  решать аналитически.
\end{enumerate}

\section{ЗАДАЧИ}\label{ux437ux430ux434ux430ux447ux438}

\begin{enumerate}
\def\labelenumi{\arabic{enumi}.}
\tightlist
\item
  Построить блок-схему игры «Блуждающий пьяница» для решения задачи
\end{enumerate}

(УЧП)

, 0 \textless{} x \textless{} 1 , 0 \textless{} y \textless{} 1 , (ГУ)
\texttt{u(x,\ y)} = \texttt{g(x,\ y)} на границе,

\begin{itemize}
\item
  во внутренних узлах сетки, если число узлов сетки произвольно. 2.
  Напишите программу для решения задачи по блок-схеме из задачи 1.
\item
  \begin{enumerate}
  \def\labelenumi{\arabic{enumi}.}
  \setcounter{enumi}{2}
  \tightlist
  \item
    Как модифицировать игру, чтобы решить задачу
  \end{enumerate}
\end{itemize}

(УЧП)

содержащей только начальные условия. Эту задачу можно было бы решить
либо с помощью преобразования Фурье (по переменной t), либо с помощью
преобразования Лапласа (по переменной t). Однако мы воспользуемся другим
методом (методом канонических координат), который, как мы надеемся,
заинтересует читателя. Этот метод базируется на тех же идеях, что и
метод перехода к движущейся системе отсчета, с которым мы познакомились
в лекции 15. Итак, приступим к решению задачи (17.1).

\section{Решение одномерного волнового уравнения. Формула
Даламбера}\label{ux440ux435ux448ux435ux43dux438ux435-ux43eux434ux43dux43eux43cux435ux440ux43dux43eux433ux43e-ux432ux43eux43bux43dux43eux432ux43eux433ux43e-ux443ux440ux430ux432ux43dux435ux43dux438ux44f.-ux444ux43eux440ux43cux443ux43bux430-ux434ux430ux43bux430ux43cux431ux435ux440ux430}

Решение задачи (17.1) разобьем на несколько шагов:

ШАГ 1 (Замена координат (x, t) новыми каноническими координатами
\texttt{(\textbackslash{}xi,\ \textbackslash{}eta)} ).

Для решения задачи (17.1) воспользуемся тем, что если заменить две
независимые переменные x и t новыми пространственновременными
координатами

, 0 \textless{} x \textless{} 1 , 0 \textless{} y \textless{} 1 , (ГУ)
\texttt{u(x,\ y)} = \texttt{g(x,\ y)} на границе?

Как выглядит в этой игре случайная прогулка?

\begin{enumerate}
\def\labelenumi{\arabic{enumi}.}
\setcounter{enumi}{3}
\tightlist
\item
  Можете ли вы так модифицировать игру, чтобы решить задачу
\end{enumerate}

\begin{verbatim}
\begin{array}{l} u_{xx} + u_{yy} = 0, \quad 0 < x < 1, \quad 0 < y < 1, \\ u(x, 1) = 0, \quad 0 < x < 1, \quad 0 < y < 1, \\ u(x, 0) = 0, \quad 0 < x < 1, \quad 0 < y < 1, \\ u(0, y) = 1, \quad \frac{\partial u}{\partial x}(1, y) = 0 \end{array}
\end{verbatim}

то уравнение

\begin{verbatim}
\begin{array}{ll} (\text{УЧП}) & u_t = \alpha^2 u_{xx} - \beta u_x - \gamma u, \quad 0 < x < 1, \quad 0 < t < \infty \\ (\text{ГУ}) & \left\{ \begin{array}{ll} u\left(0,\ t\right) = f\left(t\right), & 0 < t < \infty, \\ u\left(1,\ t\right) = g\left(t\right), & 0 < t < \infty, \\ u\left(x,\ 0\right) = \varphi\left(x\right), & 0 \leqslant x \leqslant 1. \end{array}
\end{verbatim}

преобразуется в уравнение

ПРИМІТКА. Замінимо це рівняння на скінченну різницю Кренка--Ніколсона,
виразимо \texttt{u\_\{i+1,j\}} через п'ять значень
\texttt{u\_\{i+1,j-1\}} , \texttt{u\_\{i+1,j+1\}} ,
\texttt{u\_\{i,j+1\}} , \texttt{u\_\{i,j+1\}} , \texttt{u\_\{i,j+1\}} у
суміжних точках.

\section{КАЛЬКУЛЮС ВАРІАЦІЙ (РІВНЯННЯ
ЕЙЛЕРА--ЛАГРАНЖА)}\label{ux43aux430ux43bux44cux43aux443ux43bux44eux441-ux432ux430ux440ux456ux430ux446ux456ux439-ux440ux456ux432ux43dux44fux43dux43dux44f-ux435ux439ux43bux435ux440ux430ux43bux430ux433ux440ux430ux43dux436ux430}

МЕТА ЛЕКЦІЇ: Ввести поняття функціоналу (функцію з функції) і пояснити,
як функціонали природно виникають у фізиці. Типовий функціонал є
інтегралом

.

В этом легко убедиться, если воспользоваться формулами преобразования
координат:

(17.2)

\texttt{що\ є\ функцією\ функції\ y\ (припускається,\ що\ інтегративна\ функція\ \textquotesingle{}F\textbackslash{}left(x,\textbackslash{},y,\textbackslash{},y\textquotesingle{}\textbackslash{}right)\textquotesingle{}\ задана).\ Приклад\ функціональності:}

Подставляя эти выражения для производных в волновое уравнение, получим

Ми покажемо, як знайти функцію \texttt{\textbackslash{}overline\{y\}(x)}
, яка мінімізує функціональність J{[}y{]}. Виявляється, що мінімізуюча
функція має \texttt{\textbackslash{}overline\{y\}} задовольняти так
зване рівняння Ейлера-Лагранжа. Це рівняння виконує ту ж роль, що й
необхідна умова мінімуму

.. Первый шаг сделан.

ШАГ 2 (Решение преобразованного уравнения).

Преобразованное уравнение легко решить двумя последовательными
интегрированиями (сначала по переменной ξ, а затем по η). После
интегрирования по ξ получаем

функції \texttt{f(x)} у точці x у диференціальному численні.

Варіаційне числення тісно пов'язане з диференціальними рівняннями, але,
на жаль, більшість студентів його не вивчає. Ця лекція (і наступна) є
вступом до варіаційного числення та показує, як можна розв'язувати
диференціальні рівняння в похідних на основі варіаційних принципів.

Калькулюс варіацій виник одночасно з математичним аналізом у зв'язку з
розв'язанням задач максимізації та мінімізації функцій функцій (які
називаються функціоналами). Першою задачею варіаційного числення була
задача брахістохрона, сформульована Йоганном Бернуллі у 1696 році.

Бернуллі показав, що час сходження записується у формі

\begin{verbatim}
— произвольная функция от   \eta  .
\end{verbatim}

Интегрирование последнего соотношения по \textbf{η} приводит к общему
решению

і, отже, є функцією \texttt{T\textbackslash{}\{y\textbackslash{}\}} з
функції. Оскільки багато функціоналів впорядковані таким чином, ми
зосередимося на вивченні загальної функціональності

(44.1)

\begin{verbatim}
где   \phi(\eta)   — первообразная функция   \phi_1(\eta)  , а   \psi(\xi)   — произвольная функция. Итак, общее решение уравнения
\end{verbatim}

Тепер сформульуємо головну мету наших лекцій: знайти функцію, яка дає
мінімум (або максимум) функціоналу (44.1). Стратегія пошуку буде такою
ж, як при пошуку мінімуму функції в диференціальному численні. Там ми
знайшли критичні точки функції за умовою f'(x) = 0. У варіаційному
численні все буде трохи складніше, оскільки аргумент більше не є
числовою змінною, а функцією. Однак загальний підхід залишається
незмінним: ми обчислимо так звану функціональну похідну функції
\texttt{y(x)} і прирівняємо її до нуля. Це нове рівняння буде
аналогічним

записывается в виде

(17.3)

з диференціального числення, але тепер це буде звичайне диференціальне
рівняння, відоме як рівняння Ейлера-Лагранжа. Нам просто потрібно знайти
це рівняння і розв'язати з ним деякі проблеми.

Мінімізація функціональності

\begin{verbatim}
где   \phi(\eta)   и   \psi(\xi)   — произвольные функции своих аргументов. Папример, читатель может проверить, что функции
\end{verbatim}

Розглянемо задачу пошуку функції \texttt{y(x)}, яка мінімізує
функціональність

\begin{verbatim}
u(\xi, \eta) = 1/\eta + tg \xi,    u(\xi, \eta) = \eta^2 + e
\end{verbatim}

\begin{verbatim}
являются решениями уравнения   u_{\mathrm{ }nh}=0  . Второй шаг сделан.
\end{verbatim}

\begin{verbatim}
\text{IIIA}\Gamma   3 (Возвращение к старым координатам x и t).
\end{verbatim}

Для нахождения общего решения, т. е. всех решений уравнения
\texttt{u\_\{tt\}\ =\ c\^{}2\ u\_\{xx\}},подставим

у класі гладких функцій, які задовольняють граничні умови

, \texttt{y(b)} = B (див. рис. 44.2).

\begin{verbatim}
Нехай існує y = \overline{y}(x) потрібна функція, і розглянемо малу варіацію функції \overline{y} , тобто функцію \overline{y} + \varepsilon \eta(x) , де \varepsilon — мале число, а \eta(x) — гладка функція, яка задовольняє
\end{verbatim}

\begin{verbatim}
граничні умови \eta(a) = \eta(b) = 0 (див. рис. 44.2). Очевидно, що якщо обчислити інтеграл J з близької функції y + \varepsilon \eta , то функціонал буде збільшуватися, тобто
\end{verbatim}

в решение

\begin{verbatim}
за всі \epsilon . Іншими словами, графік \phi\left(\epsilon\right)=J\left[\overline{y}+\epsilon\eta\right] як функція, \epsilon буде виглядати як рис. 44.3.
\end{verbatim}

\pandocbounded{\includegraphics[keepaspectratio]{images/_page_329_Figure_5.jpeg}}

Рис. 44.1. Задача брахістохрана (варіаційне числення починалося з цієї
задачі).

\pandocbounded{\includegraphics[keepaspectratio]{images/_page_329_Figure_7.jpeg}}

\begin{verbatim}
Рис. 44.2. Варіація функції: a — варіація функції \overline{y} = \overline{y}(x) + \epsilon \eta(x) ; \delta є мінімізуючою функцією \overline{y}(x) .
\end{verbatim}

З Фіг. 44.3 Очевидно, що в цьому випадку необхідно обчислити похідну від

В результате получаем

(17.4)

\begin{verbatim}
Po E, поставити \varepsilon = 0 і прирівняти отриманий вираз до нуля, тобто
\end{verbatim}

Это общее решение волнового уравнения. С физической точки эрения оно
интересно тем, что представляет сумму двух бегущих волн Р И С

\pandocbounded{\includegraphics[keepaspectratio]{images/_page_126_Figure_14.jpeg}}

. 17.1. Волна \texttt{u(x,\ t)\ =\ e\^{}\{-(x-ct)\^{}2\}} движется слева
направо

произвольной формы, движущихся в прогивоположных направлениях со
скоростью c.~Например, функции

\begin{verbatim}
= \int_{0}^{d} \left[\frac{\partial F}{\partial \overline{y}} \eta\left(x\right) + \frac{\partial F}{\partial \overline{y'}} \eta'\left(x\right)\right] dx = 0.
\end{verbatim}

(волна, Лвижущаяся слева направо),

\begin{verbatim}
\frac{d\varphi(e)}{de} \equiv \int_{a}^{b} \left\{ \frac{\partial F}{\partial \overline{y}} - \frac{\partial}{\partial x} \left[ \frac{\partial F}{\partial \overline{y'}} \right] \right\} \, \eta(x) \, dx = 0.
\end{verbatim}

(волна, движущаяся справа налево),

\begin{verbatim}
\frac{\partial F}{\partial \overline{y}} - \frac{d}{dx} \left[ \frac{\partial F}{\partial \overline{y'}} \right] = 0
\end{verbatim}

(две волны, движущиеся в противоположных направлениях)

являются типичными решениями волиового уравнения. На рисунке 17.1
изображена простейшая бегущая волна.

ШАГ 4 (Учет начальных условий).

Напомним читателю общий метод решения задачи Коши в теории обыкновенных
дифференциальных уравнений. Сначала находят общее решение уравнения, а
затем оно подставляется в начальные условия для того, чтобы найти
конкретные значения произвольных постоянных. У нас ситуация аналогична.
Для решения задачи Коши для волнового уравнения мы подставим общее
решение

\begin{verbatim}
\frac{\partial F}{\partial y} - \frac{d}{dx} \left[ \frac{\partial F}{\partial y'} \right] = 0
\end{verbatim}

(содержащее две произвольные функции) в начальные условия

\begin{verbatim}
J[y] = \int_{0}^{1} [y^{2} + y'^{2}] dx.
\end{verbatim}

\texttt{u\_1(x,\ 0)\ =\ g(x),} чтобы найти конкретные выражения для
произвольных функций ф и ф. После подстановки получаем

(17.5)

\begin{verbatim}
F_y - \frac{d}{dx} F_{y'} = 0,
\end{verbatim}

Если теперь проинтегрировать второе уравнение (17.5), то получится
алгебранческая система двух уравнений относительно неизвестных функций
\texttt{\textbackslash{}varphi(x)} и \texttt{\textbackslash{}psi(x)} . В
самом деле, после интегрирования в пределах от \texttt{x\_0} до x
получаем

(17.6)

\begin{verbatim}
\begin{aligned} F_y &= 2y, \ F_{y'} &= 2y'. \end{aligned}
\end{verbatim}

Решаем теперь уравнение (17.6) совместно с первым уравнением (17.5) и
получаем следующие выражения для функций
\texttt{\textbackslash{}varphi(x)} и \texttt{\textbackslash{}psi(x)} :

(17.7)

\begin{verbatim}
2y - \frac{d}{dx}(2y') = 0,
y'' - y = 0.
\end{verbatim}

Следовательно, общее решение задачи (17.1) дается формулой

(17.8)

\begin{verbatim}
\overline{y}(x) = 0.42e^x - 0.42e^{-x}
\end{verbatim}

Цель достигнута. Мы получили общее решение вадачи (17.1). Решение вида
(17.8) принято называть формулой Даламбера. Читателю предлагается
самостоятельно проследить за расстановкой пределов интегрирования, чтобы
убедиться в том, что они меняются от x-ct до x+ct. Задача полностью
решена.

Прежде чем завершить лекцию, приведем несколько примеров, показывающих,
как применять формулу Даламбера в конкретных валачах.

\section{Примеры применения формулы
Даламбера}\label{ux43fux440ux438ux43cux435ux440ux44b-ux43fux440ux438ux43cux435ux43dux435ux43dux438ux44f-ux444ux43eux440ux43cux443ux43bux44b-ux434ux430ux43bux430ux43cux431ux435ux440ux430}

\section{1. Движение синусондальной
волны}\label{ux434ux432ux438ux436ux435ux43dux438ux435-ux441ux438ux43dux443ux441ux43eux43dux434ux430ux43bux44cux43dux43eux439-ux432ux43eux43bux43dux44b}

Рассмотрим начальные условия вида

\begin{verbatim}
J[y] = \int_a^b y \sqrt{1 + y'^2} \, dx,
\end{verbatim}

\texttt{u\_1}(x, 0) = 0. По формуле Даламбера получаем решение

\begin{verbatim}
\overline{y}(x) = \alpha \operatorname{ch}[(x-\beta)/\alpha],
\end{verbatim}

\begin{verbatim}
Его можно интерпретировать следующим образом. Начальное смещение струны   u(x, 0) = \sin x   делится на две одинаковые части
\end{verbatim}

\begin{verbatim}
\overline{y}(x) = 0.42e^x - 0.42e^{-x}
H   \frac{\sin x}{2}  .
\end{verbatim}

Каждая из частей распространяется со скоростью с в виде бегущих волн.
Одна из волн перемещается слева направо, а вторая---в противоположном
направлении. Читателю предлагается подумать над тем, как выглядит в этом
случае результирующая волна.

\section{2. Распространение простейшего прямоугольного
импульса}\label{ux440ux430ux441ux43fux440ux43eux441ux442ux440ux430ux43dux435ux43dux438ux435-ux43fux440ux43eux441ux442ux435ux439ux448ux435ux433ux43e-ux43fux440ux44fux43cux43eux443ux433ux43eux43bux44cux43dux43eux433ux43e-ux438ux43cux43fux443ux43bux44cux441ux430}

В этом случае начальные условия задаются следующим образом:

\begin{verbatim}
J[y] = \int_{0}^{1} [y^{2} + y'^{2}] dx.
\end{verbatim}

\texttt{u\_t}(x, 0) = 0 (в остальных точках).

Поскольку начальное смещение расшепляется на две полуволны, движущиеся в
противоположных направлениях, результирующее движение будет таким, как
показано на рис. 17.2.

\pandocbounded{\includegraphics[keepaspectratio]{images/_page_129_Figure_2.jpeg}}

РИС, 17.2. Начальное возмущение распадается на две бегущие волны.

\section{3. Задана начальная
скорость}\label{ux437ux430ux434ux430ux43dux430-ux43dux430ux447ux430ux43bux44cux43dux430ux44f-ux441ux43aux43eux440ux43eux441ux442ux44c}

\begin{verbatim}
Предположим теперь, что в начальный момент струна находится в положении равновесия. Придадим струне начальную скорость (как в фортепиано) вида   \sin x  , т. е.
\end{verbatim}

\begin{verbatim}
\int\limits_{t_1}^{t_2}
J\left[u\right] = \iint\limits_{D} F\left(x,\ y,\ u,\ u_{x},\ u_{y}\right) \, dx \, dy.
\end{verbatim}

Решение находим по формуле Даламбера

\begin{verbatim}
F_{u} - \frac{\partial}{\partial x} F_{ux} - \frac{\partial}{\partial y} F_{uy} = 0.
\end{verbatim}

Оно представляет сумму двух бегущих косинусоидальных воли. Читатель
должен спросить себя, кажется ли ему это решение разумным.

На этом лекция 17 заканчивается. В следующей лекции мы покажем, как
формула Даламбера помогает дать интерпретацию в плоскости переменных x и
t.

\section{ЗАМЕЧАНИЯ}\label{ux437ux430ux43cux435ux447ux430ux43dux438ux44f-1}

\begin{itemize}
\tightlist
\item
  \begin{enumerate}
  \def\labelenumi{\arabic{enumi}.}
  \tightlist
  \item
    Общее решение уравнения с частными производными второго порядка
    содержит две произвольные функции, а общее решение обыкновенного
    дифференциального уравнения содержит две произвольные постоянные.
    Значит, у уравнения с частными производными больше решений, чем у
    обыкновенного дифференциального уравнения.
  \end{enumerate}
\item
  2 Замена переменных для приведения уравнения с частными производными к
  более простому виду является одним из общих методов теории. Новые
  координаты (ξ, η) принято называть каноническими координатами. Мы
  будем обращаться к ним и в дальнейшем, особенно при изучении
  гиперболических уравнений.
\item
  3 Метод нахождения общего решения уравнения с частными производными и
  подстановки этого решения в начальные условия не является общим
  методом решения уравнений с частными производными. Решение, которое
  рассматривается в данной лекции,--- единственный пример использования
  этого метода. Обычно невозможно найти общее решение уравнения с
  частными производными, но даже если удается это сделать, оказывается
  очень сложным подставить это решение в краевые условия.
\end{itemize}

\section{ЗАДАЧН}\label{ux437ux430ux434ux430ux447ux43d}

\begin{enumerate}
\def\labelenumi{\arabic{enumi}.}
\item
  Проверьте, что решение Даламбера (17-8) удовлетворяет начальным
  условиям задачи (17.1).
\item
  Подставьте (17.7) в общее решение
  \texttt{u(x,\ t)\ =\ \textbackslash{}psi(x\ -\ ct)\ +\ \textbackslash{}psi(x\ +\ \textbackslash{}epsilon\ t)}
  и получите формулу Даламбера.
\item
  Найдите решение задачи Коши:
\end{enumerate}

\begin{verbatim}
J[y] = \int_{0}^{1} \sqrt{1 + y'^{2}} dx.
\end{verbatim}

Как выглядит это решение в различные моменты времени? 4. Найдите решение
задачи Коши:

(УЧП)

\begin{verbatim}
\int_{t_{1}}^{t_{2}} [KE - PE] dt = \frac{1}{2} \int_{t_{1}}^{t_{2}} [m\dot{y}^{2} - ky^{2}] dt
\end{verbatim}

(НУ)

\begin{verbatim}
J[y] = \int_{0}^{\pi/2} [y'^{2} - y^{2}] dx
\end{verbatim}

Постройте график решения в различные моменты времени. 5. Решая совместно
уравнения (17.6) и первое уравнение из (17.5),

получите формулы (17.7).

\section{ФОРМУЛА ДАЛАМБЕРА
(ПРОДОЛЖЕНИЕ)}\label{ux444ux43eux440ux43cux443ux43bux430-ux434ux430ux43bux430ux43cux431ux435ux440ux430-ux43fux440ux43eux434ux43eux43bux436ux435ux43dux438ux435}

ЦЕЛЬ ЛЕКЦИИ! Показать, как с помощью формулы Даламбера можно найти
решение волнового уравнения на полуограниченной прямой

(УЧП)

\begin{verbatim}
J[u] = \iint_D F[x, y, u, u_x, u_y) \, dx \, dy,
\end{verbatim}

(ГУ)
\texttt{u(0,\ t)\ =\ 0,\ \textbackslash{}quad\ 0\ \textless{}\ t\ \textless{}\ \textbackslash{}infty,}
(НУ)

\begin{verbatim}
F_{u} - \frac{\partial}{\partial x} F_{uv} - \frac{\partial}{\partial y} F_{uv} = 0.
\end{verbatim}

Дополнительно дается интерпретация формулы Даламбера в координатной
плоскости x, t.

На прошлой лекции было показано, что если заданы начальное смещение
струны \texttt{u(x,\ 0)} = \texttt{f(x)} и начальная скорость
\texttt{u\_t(x,\ 0)\ =\ g(x)} , то выражение

(18.i)

\begin{verbatim}
u_{xx} + u_{yy} = f
\end{verbatim}

описывает смещение струны в произвольный момент времени. В этой лекции
мы познакомим читателя с одной интересной интерпретацией этой формулы в
плоскости переменных x и t (пространственновременная или фазовая
плоскость) и покажем, как изменить формулу Даламбера, чтобы получить
решение задачи о колебаниях полубесконечной струны.

Начнем с интерпретации формулы (18.1) в плоскости переменных x и t.

\section{Пространственно-временная интерпретация формулы
Даламбера}\label{ux43fux440ux43eux441ux442ux440ux430ux43dux441ux442ux432ux435ux43dux43dux43e-ux432ux440ux435ux43cux435ux43dux43dux430ux44f-ux438ux43dux442ux435ux440ux43fux440ux435ux442ux430ux446ux438ux44f-ux444ux43eux440ux43cux443ux43bux44b-ux434ux430ux43bux430ux43cux431ux435ux440ux430}

На предыдущей лекции мы показали, что решение задачи Коши

(18.2)

\begin{verbatim}
J[u] = \iint_{D} [u_{x}^{2} + u_{y}^{2} + 2uf] dx dy.
J[u] = \int_{0}^{1} \int_{0}^{1} \left[ u_{1}^{2} + u_{y}^{2} \right] dx dy
u(x, t) = \frac{1}{2} [f(x-ct) + f(x+ct)] + \frac{1}{2c} \int_{x-ct}^{x+ct} g(\xi) d\xi.
u_{xx} + u_{yy} = 0
u(x, 0) = f(x),
u_{xx} + u_{yy} = 0
u(x, t) = \frac{1}{2} [f(x-ct) + f(x+ct)].
J[u] = \int_{0}^{1} \int_{0}^{1} \left[ u_{x}^{2} + u_{y}^{2} \right] dx dy
x-ct=x_0-ct_0
\end{verbatim}

(PPP)

\begin{verbatim}
(Y \text{ ЧП}) \qquad u_{tt} = c^2 u_{xx}, \quad -\infty < x < \infty, \quad 0 < t < \infty,
J[u] = \iint_{D} [u] + u_{\eta} + 2uf ] dx dy.
(HY) \qquad
u_{xx} + u_{yy} = f
J[u] = \int_{0}^{1} \int_{0}^{1} [u' + u' + 2uf] dx dy,
J[u] = \int_{0}^{1} \int_{0}^{1} \left[ u_{x}^{2} + u_{y}^{2} + 2uf \right] dx dy,
\end{verbatim}

На рис. 18.2 решение изображено в плоскости переменных x, t, а не u, x,
как в предыдущей лекции.

\pandocbounded{\includegraphics[keepaspectratio]{images/_page_133_Figure_5.jpeg}}

РИС. 18.2. Решение задачи Коши (18.2) в xt-плоскости a --- передний
фронт волны, движущейся налево; b --- задний фронт волны; b --- переднии
фронт волны, движущейся направо

Теперь займемся интерпретацией формулы Даламбера, когда начальное
смещение равно нулю, а начальная скорость произвольна.

СЛУЧАЙ 2 (Начальное смещение равно нулю, начальная скорость
произвольна).

Рассмотрим начальные условия вида

\begin{verbatim}
u_n(x, y) = a_1 \varphi_1(x, y) + a_2 \varphi_2(x, y) + \ldots + a_n \varphi_n(x, y),
\end{verbatim}

\texttt{u\_t}(x, 0) = \texttt{g(x)}. Решение в этом случае имеет вид

\begin{verbatim}
\varphi_1(x, y) = xy(1-x)(1-y)
\end{verbatim}

\begin{verbatim}
Значение величины u в точке   (x_{\mathrm{ }e},\ t_{\mathrm{ }e})   можно интерпретировать как интеграл от начальной скорости в пределах от   x_{\mathrm{ }e}-ct_{\mathrm{ }e}   до Р И С
\end{verbatim}

\pandocbounded{\includegraphics[keepaspectratio]{images/_page_134_Figure_2.jpeg}}

. 18.3. Интерпретация задачи с начальной скоростью в x1-плоскости.

\begin{verbatim}
x_{\bullet}+ct_{\bullet}   (рис. 18.3). В качестве примера на рис. 18.4 приводится графическая интерпретация решения задачи
\end{verbatim}

(13.4)

\begin{verbatim}
u_{1}(x, y) = a_{1}xy(1-x)(1-y),
u_{3}(x, y) = xy(1-x)(1-y)[a_{1} + a_{2}x + a_{3}y],
\vdots
u(x, t) = \frac{1}{2c} \int_{x-ct}^{x+ct} g(\xi) d\xi =
J[u_n] = \int_0^1 \int_0^1 \left\{ \left[ \sum_{j=1}^n a_j \frac{\partial \varphi_j}{\partial x} \right]^2 + \left[ \sum_{j=1}^n a_j \frac{\partial \varphi_j}{\partial y} \right]^2 + 2f \sum_{j=1}^n a_j \varphi_j \right\} dx dy
= \frac{1}{2c} \int_{x-ct}^{x+ct} 0 d\xi, \quad (x, t) \in \text{ области 1,}
\frac{\partial J\left[u_{n}\right]}{\partial a_{1}} = 2 \int_{0}^{1} \int_{0}^{1} \left\{ \sum_{j=1}^{n} \left[ \frac{\partial \varphi_{j}}{\partial x} \frac{\partial \varphi_{j}}{\partial x} + \frac{\partial \varphi_{j}}{\partial y} \frac{\partial \varphi_{1}}{\partial y} \frac{\partial \}{\partial y} \right] a_{j} + f \varphi_{1} \right\} dx dy = 0,
= \frac{1}{2c} \int_{x-ct}^{x+ct} d\xi, \quad (x, t) \in \text{ области 2,}
\end{verbatim}

Все это представляется изрядно сложным, но если переписать все эти
уравнения в матричной форме, то получится система линейных уравнений

\begin{verbatim}
= \frac{1}{2c} \int_{x-ct}^{-1} d\xi, \quad (x, t) \in \text{ области 3,}
\end{verbatim}

где \texttt{A\ =\ (A\_\{IJ\})} матрица размером
\texttt{n\ \textbackslash{}times\ n} , элементы которой вычисляются по
формулам

(45.3)

\begin{verbatim}
= \frac{1}{2c} \int_{x-ct}^{1} d\xi, \quad (x, t) \in \text{ области 4,}
\end{verbatim}

\texttt{b\ =\ (b\_t)} --- вектор с компонентами

\begin{verbatim}
= \frac{1}{2c} \int_{x-ct}^{x+ct} 0 d\xi, \quad (x, t) \in \text{ области 5,}
\end{verbatim}

\texttt{а\ \ \textquotesingle{}a=(a\_i)\textquotesingle{}\ —\ неизвестный\ вектор,\ компоненты\ которого\ являются\ коэффициентами\ в\ приближенном\ решении\ \ \textquotesingle{}a\_n(x,y)=a\_1\textbackslash{}varphi\_1(x,y)+a\_2\textbackslash{}varphi\_2(x,y)+\textbackslash{}ldots+a\_n\textbackslash{}varphi\_n(x,y)\textquotesingle{}\ .}

ШАГ 3. Решаем линейную систему Aa = b относительно коэффициентов
\texttt{a\_1,\ a\_2,\ \textbackslash{}ldots,\ a\_n} и получаем
приближенную минимизирующую функцию

\begin{verbatim}
= \frac{1}{2c} \int_{x-ct}^{x+ct} d\xi, \quad (x, t) \in \text{ области 6.}
\end{verbatim}

а эначит, и приближенное решение задачи Дирихле.

\section{ЗАМЕЧАНИЯ}\label{ux437ux430ux43cux435ux447ux430ux43dux438ux44f-2}

\begin{enumerate}
\def\labelenumi{\arabic{enumi}.}
\tightlist
\item
  В этой лекции методом Ритца минимизировался функционал, вависящий от
  функции двух переменных. Точно так же можно минимизировать функционалы
  вида
\end{enumerate}

\begin{verbatim}
y(0) = 0, \quad y(1) = 1.
\begin{array}{llll} (\text{Y}4\Pi) & u_{tt} = c^2 u_{xx}, & 0 < x < \infty, & 0 < t < \infty, \\ (\Gamma \text{Y}) & u(0, t) = 0, & 0 < t < \infty, \\ (\text{H}\text{Y}) & \begin{cases} u(x, 0) = f(x), \\ u_t(x, 0) = g(x). \end{cases} & 0 < x < \infty, \\ \end{array}
\varphi_i(0) = 0,
J[u] = \int_{0}^{1} \int_{0}^{1} [u^{2} + u^{2}_{y}] dx dy,
\end{verbatim}

\section{Решение задачи для полубесконечной струны с помощью формулы
Даламбера}\label{ux440ux435ux448ux435ux43dux438ux435-ux437ux430ux434ux430ux447ux438-ux434ux43bux44f-ux43fux43eux43bux443ux431ux435ux441ux43aux43eux43dux435ux447ux43dux43eux439-ux441ux442ux440ux443ux43dux44b-ux441-ux43fux43eux43cux43eux449ux44cux44e-ux444ux43eux440ux43cux443ux43bux44b-ux434ux430ux43bux430ux43cux431ux435ux440ux430}

Наша задача---найти волновые движения полубесконечной струны, левый
конец которой закреплен, при заданных начальных условиях. Для решения
задачи (18.6) мы воспользуемся тем же приемом, что и при отыскании
общего решения волнового уравнения на всей прямой. Подставляя общее
решение

\begin{verbatim}
u_{xx} + u_{yy} = 1
\end{verbatim}

в начальные условия, находим (как в предыд, щей лекции)

(18.7)

\begin{verbatim}
J[y] = \int_{0}^{1} [y^{2} + {y'}^{2}] dx, \ y(0) = 0, \ y(1) = 1,
z(x) = (1 - x) y(x),
\end{verbatim}

\texttt{"Теперь\ перед\ нами\ проблема,\ которой\ не\ было\ в\ задаче\ о\ бесконечной\ струне.\ Решение\ \ \textquotesingle{}u\textbackslash{}left(x,\textbackslash{}t\textbackslash{}right)\textquotesingle{}\ \ должно\ быть\ определено\ всюду\ только\ внутри\ первого\ квадранта\ \ \textquotesingle{}(x\textgreater{}0,\textbackslash{},t\textgreater{}0)\textquotesingle{}\ \ плоскости\ переменных\ x\ и\ t.\ Значит,\ мы\ должны\ уметь\ вычислить\ функции}

\begin{verbatim}
u_{xx} + u_{yy} = 0
\end{verbatim}

при всех
\texttt{-\textbackslash{}infty\ \textless{}\ x-ct\ \textless{}\ \textbackslash{}infty}
, \texttt{\textbackslash{}psi(x+ct)} при всех
\texttt{0\ \textless{}\ x+ct\ \textless{}\ \textbackslash{}infty} .

\texttt{К\ сожалению,\ первая\ из\ формул\ (18.7)\ позволяет\ вычислять\ \ \textquotesingle{}\textbackslash{}varphi(x-ct)\textquotesingle{}\ \ только\ при\ \ \textquotesingle{}x-ct\textbackslash{}geqslant\ 0\textquotesingle{}\ ,\ поскольку\ в\ начальных\ данных\ функции\ \textquotesingle{}f(x)\textquotesingle{}\ и\ \textquotesingle{}g(x)\textquotesingle{}\ определены\ только\ для\ положительных\ значений\ аргумента.}

Если \texttt{x-ct\ \textbackslash{}geqslant\ 0} , то после подстановки
(18.7) в общее решение

\begin{verbatim}
J[u] = \int_{0}^{1} \int_{0}^{1} [u_{x}^{2} + u_{y}^{2}] dx dx.
\end{verbatim}

получаем

\begin{verbatim}
u_{xx} + u_{yy} = \sin(\pi x)
u(x, y) = \left[Ae^{\pi y} + Be^{-\pi y} - \frac{1}{n^2}\right] \sin(\pi x),
\begin{split} u\left(x,\,t\right) &= \phi\left(x-ct\right) + \psi\left(x+ct\right) = \\ &= \frac{1}{2} \left[ f\left(x-ct\right) + f\left(x+ct\right) \right] + \frac{1}{2c} \int\limits_{x-ct}^{x+ct} g\left(\xi\right) \, d\xi. \end{split}
\end{verbatim}

\texttt{\textbackslash{}Дельта\ u\ +\ u\^{}2\ =\ 0}

\begin{verbatim}
\Delta u = 0
\end{verbatim}

Что же делать, если x \textless{} ct? Пришла пора воспользоваться
граничным условием \texttt{u(0,t)}=0. С помощью этого граничного условия
мы доопределяем функцию \texttt{\textbackslash{}varphi(x-ct)} при x
\textless{} ct. Для этого подставим общее решение
\texttt{u(x,t)=\textbackslash{}varphi(x-ct)+\textbackslash{}psi(x+ct)} в
граничное условие \texttt{u(0,t)}=0. В результате получим
\texttt{\textbackslash{}varphi(-ct)=-\textbackslash{}psi(ct)} ,

откуда

\begin{verbatim}
\Delta u = 0
\end{verbatim}

Подстановка этого выражения для ф в общее решение дает

\begin{verbatim}
\Delta u + \varepsilon u^2 = 0
\end{verbatim}

Комбинируя решения для x \textless{} ct и x \textgreater{} ct,
окончательно получаем (см. рис. 18.6)

\texttt{\textbackslash{}Дельта\ u\ +\ u\^{}2\ =\ 0}

\begin{verbatim}
\Delta u = 0 \rightarrow \Delta u + \varepsilon u^2 = 0 \rightarrow \Delta u + u^2 = 0
\begin{cases} \frac{1}{2} \left[ f(x - ct) + f(x + ct) \right] + \frac{1}{2c} \int_{x - ct}^{x + ct} g(\xi) d\xi, & x \ge ct, \\ \frac{1}{2} \left[ f(x + ct) - f(ct - x) \right] + \frac{1}{2c} \int_{ct - x}^{x + ct} g(\xi) d\xi, & x < ct. \end{cases}
u_0 \rightarrow u = \sum_{k=0}^{\infty} \varepsilon^k u_k \rightarrow u = \sum_{k=0}^{\infty} u_k
\Delta u + cu^2 = 0,
\end{verbatim}

Этим завершается наша лекция. Интерпретация решения (18.8) будет
приведена в замечаниях.

\pandocbounded{\includegraphics[keepaspectratio]{images/_page_137_Figure_9.jpeg}}

РИС. 18.6. Интерпретация задачи Коши для полуограниченной струны в
xt-плоскости a---возмущение отражается от границы;
\texttt{\textbackslash{}delta} ---возмущение распространяется вдоль
характеристики из точки ( \texttt{x\_0+ct\_0} , 0).

\section{ЗАМЕЧАНИЯ}\label{ux437ux430ux43cux435ux447ux430ux43dux438ux44f-3}

\begin{enumerate}
\def\labelenumi{\arabic{enumi}.}
\item
  Решение (18.8) полностью соогветствует нашим представлениям о
  полубесконечной струне с граничным условием \texttt{u(0,\ t)} = 0. При
  \texttt{x\ \textbackslash{}geqslant\ ct} оно совпадает с формулой
  Даламбера для бесконечной струны. Если x \textless{} ct, то формула
  модифицируется таким образом, что в нее включается волна, отраженная
  от границы (при отражении знак волны меняется на противоположный).
\item
  Если граничное условие \texttt{u(0,\ t)} = 0 изменить на другое, то и
  решение (18.8), естественно, изменится. Можно получить в явном виде
  решения для граничных условий
\end{enumerate}

\begin{verbatim}
\Delta u + \varepsilon u^2 = 0
\end{verbatim}

или

\begin{verbatim}
(46.1) u = u_0 + \varepsilon u_1 + \varepsilon^2 u_2 + \dots,
\end{verbatim}

Дополнительную информацию по этим вопросам читатель может найти в
литературе, указанной в конце книги.

\begin{enumerate}
\def\labelenumi{\arabic{enumi}.}
\setcounter{enumi}{2}
\tightlist
\item
  Прямые линии
\end{enumerate}

\texttt{\textbackslash{}Дельта\ u\ +\ u\^{}2\ =\ 0}

\texttt{x\ -\ ct\ =\ \textbackslash{}text\{const\}} называются
характеристиками. Возмущение распространяется вдоль характеристик.
Понятие характеристики тесно связано с уравнениями гиперболического
типа.

\section{ЗАДАЧИ}\label{ux437ux430ux434ux430ux447ux438-1}

\begin{enumerate}
\def\labelenumi{\arabic{enumi}.}
\tightlist
\item
  Решите смещанную задачу для полубесконечной струны
\end{enumerate}

\begin{verbatim}
\Delta u + u^2 = 0,
\Delta u = 0, \quad 0 < r < 1, \quad 0 \le \theta < 2\pi,
\end{verbatim}

Постройте графики решения для разл чных моментов времени.

\begin{enumerate}
\def\labelenumi{\arabic{enumi}.}
\setcounter{enumi}{1}
\tightlist
\item
  Для решения предыдущей задачи можно воспользоваться следующим
  алгоритмом:
\end{enumerate}

а) продолжить начальные условия на всю действительную ось по формулам

имеет только вулевое решение. - Прим. ред.

\begin{enumerate}
\def\labelenumi{\arabic{enumi})}
\setcounter{enumi}{1}
\item
  Здесь и ниже речь идет об отыскании одного из решений (второе таким
  методом не находится).--- Прим ред.
\item
  Эта «ракета», во-первых, не при всех граничных условиях долетает до
  є=1 (где при опредсленных вещественных граничных условиях открывается
  мир комплексных решений) и, во-вторых, ни при каких в не долетает до
  второго вещественного решения, когда оно существует.--- Прим. ред.
\end{enumerate}

ции \texttt{u\_0} , \texttt{u\_1} , \texttt{u\_2} , \ldots, входящие в
степенной ряд (46.1). В оставшейся части лекции мы покажем, как этим
методом можно решить различные краевые задачи.

\section{\texorpdfstring{Решение нелинейного уравнения
\texttt{\textbackslash{}Delta\ u\ +\ u\^{}2\ =\ 0} методом
возмущений}{Решение нелинейного уравнения \textbackslash Delta u + u\^{}2 = 0 методом возмущений}}\label{ux440ux435ux448ux435ux43dux438ux435-ux43dux435ux43bux438ux43dux435ux439ux43dux43eux433ux43e-ux443ux440ux430ux432ux43dux435ux43dux438ux44f-delta-u-u2-0-ux43cux435ux442ux43eux434ux43eux43c-ux432ux43eux437ux43cux443ux449ux435ux43dux438ux439}

Предположим, что нам нужно решить следующую нелинейную задачу Дирихле:

(46.2)

\begin{verbatim}
u_t(x, 0) = 0, -\infty < x < 0.  б) усреднить две волны, движущиеся налево и направо (как это было в предыдущей лекции);
\end{verbatim}

\begin{verbatim}
**в**) рассмотреть полученное решение при   x \ge 0  .
\end{verbatim}

Используя этот алгоритм, постройте графики решения для различных
моментов времени t, если начальное условие задано графически.

\pandocbounded{\includegraphics[keepaspectratio]{images/_page_139_Figure_3.jpeg}}

\begin{enumerate}
\def\labelenumi{\arabic{enumi}.}
\setcounter{enumi}{2}
\tightlist
\item
  Решите смешанную задачу для полубесконечной струны
\end{enumerate}

Розглянемо цю нелінійну задачу як збурену лінійну задачу:

(46.3)

\begin{verbatim}
Розв'язок якої u_0(r, \theta) = r \cos \theta відомий. Як змінити u_0(r, \theta) , щоб отримати рішення проблеми (46.2)? Як уже згадувалося, розглянемо клас задач \Delta u + \epsilon u^2 = 0 з граничними умовами і шукаємо розв'язок кожної задачі з цього класу у вигляді
\end{verbatim}

(46.4)

Воспользуйтесь методом, аналогичным рассмотренному в лекции. Дайте
физическую интерпретацию этой задачи.

\begin{enumerate}
\def\labelenumi{\arabic{enumi}.}
\setcounter{enumi}{3}
\tightlist
\item
  Предположим, что колебания струны описываются уравнением
  \texttt{u\_\{tt\}\ =\ u\_\{xx\}} , начальное смещение изображено на
  рисунке, а начальная скорость равна нулю. Изобразите решение этой
  задачи в xt-плоскости. Обратите внимание на то, что начальные условия
  в этой задаче разрывны.
\end{enumerate}

\pandocbounded{\includegraphics[keepaspectratio]{images/_page_139_Figure_8.jpeg}}

\section{ВОЛНОВОЕ УРАВНЕНИЕ И ГРАНИЧНЫЕ
УСЛОВИЯ}\label{ux432ux43eux43bux43dux43eux432ux43eux435-ux443ux440ux430ux432ux43dux435ux43dux438ux435-ux438-ux433ux440ux430ux43dux438ux447ux43dux44bux435-ux443ux441ux43bux43eux432ux438ux44f}

ЦЕЛЬ ЛЕКЦИИ: Познакомыть читателя с основными типами граничных условий
для волнового уравнения

\begin{verbatim}
Потім знаходимо розв'язок задачі (46.2), поклавши \varepsilon = 1^{\circ} ).
\end{verbatim}

Вираз (46.4), який претендує на розв'язок, буде замінений на обурену
проблему

(PPP)

на конечном отрезке.

В дальнейшем будут рассматриваться три основных типа граничных условий.

\begin{enumerate}
\def\labelenumi{\arabic{enumi}.}
\tightlist
\item
  Заданный режим (граничные условия 1-го рода)
\end{enumerate}

,

\begin{verbatim}
(GI) u(1, \theta) = \cos \theta .
\end{verbatim}

У результаті ми отримуємо

\texttt{u(L,\ t)\ =\ g\_2(t).} 2. Заданные силы (граничные условия 2-го
рода)

\begin{verbatim}
u_0 (1, \theta) + \varepsilon u_1 (1, \theta) + \varepsilon^2 u_2 (1, \theta) + \ldots = \cos \theta.
\end{verbatim}

\texttt{u\_x(L,\ t)\ =\ g\_2(t).} 3. Упругое закрепление (граничные
условия 3-го рода)

\begin{verbatim}
P_{0}
\end{verbatim}

\texttt{u\_x(L,\ t)\ -\ \textbackslash{}gamma\_2\ u(L,\ t)\ =\ g\_2(t).}
Могут встречаться и произвольные комбинации этих трех типоз граничных
условий. В лекции будет показан физический смысл этих граничных условий.

Итак, пока что мы рассмотрели единственный вид волнового
движения---поперечные колебания струны. Читатель должен понимать, что
это только один из видов волновых движений.

Ниже приведен перечень основных типов волн:

\begin{itemize}
\tightlist
\item
  \begin{enumerate}
  \def\labelenumi{\arabic{enumi}.}
  \tightlist
  \item
    Звуковые волны (продольные волны).
  \end{enumerate}
\item
  \begin{enumerate}
  \def\labelenumi{\arabic{enumi}.}
  \setcounter{enumi}{1}
  \tightlist
  \item
    Электромагнитные волны, в том числе световые.
  \end{enumerate}
\item
  \begin{enumerate}
  \def\labelenumi{\arabic{enumi}.}
  \setcounter{enumi}{2}
  \tightlist
  \item
    Колебания твердых тел (продольные, поперечные и крутильные).
  \end{enumerate}
\item
  \begin{enumerate}
  \def\labelenumi{\arabic{enumi}.}
  \setcounter{enumi}{3}
  \tightlist
  \item
    Волны вероятности в квантовой механике.
  \end{enumerate}
\item
  \begin{enumerate}
  \def\labelenumi{\arabic{enumi}.}
  \setcounter{enumi}{4}
  \tightlist
  \item
    Волны на воде (поперечные волны).
  \end{enumerate}
\item
  \begin{enumerate}
  \def\labelenumi{\arabic{enumi}.}
  \setcounter{enumi}{5}
  \tightlist
  \item
    Колебания струны (поперечные волны).
  \end{enumerate}
\end{itemize}

Цель нашей лекции --- обсудить различные типы граничных условий, которые
возникают при решении физических задач волнового характера. Мы
остановимся только на одномерных задачах и линейных граничных условиях.
Обычно различают граничные условия трех типов

\begin{enumerate}
\def\labelenumi{\arabic{enumi}.}
\tightlist
\item
  Заданные режимы в граничных точках (граничные условия 1-го рода)
\end{enumerate}

\begin{verbatim}
\qquad u_{\theta}(r, \theta) = r \cos \theta,
\end{verbatim}

\texttt{u(L,\ t)\ =\ g\_2(t).} 2. Заданные силы в граничных точках
(граничные условия 2-го рода)

\begin{verbatim}
\begin{cases} \Delta u_{1} = -u_{0}^{2}, \\ u_{1}(1, \theta) = 0, \\ u_{2}(1, \theta). \end{cases}
\end{verbatim}

\texttt{u\_x(L,\ t)\ =\ g\_2(t).} 3. Упругое закрепление в граничных
точках (граничные условия 3-го рода)

\begin{verbatim}
\begin{cases} \Delta u_{2} = -2u_{0}u_{1}, \\ u_{2}(1, \theta). \end{cases}
\end{verbatim}

\texttt{u\_x(L,\ t)\ -\ \textbackslash{}gamma\_2\ u(L,\ t)\ =\ g\_2(t).}
Мы начнем с граничных условий 1-го рода.

\begin{enumerate}
\def\labelenumi{\arabic{enumi}.}
\tightlist
\item
  Заданные режимы в граничных точках.
\end{enumerate}

Нам предстоит решить задачу вида

(19.1)

\begin{verbatim}
P_{1}
\end{verbatim}

Щоб розв'язати \emph{неоднорідне рівняння}, ми також звернемося до
методу, який використовувався в теорії звичайних диференціальних рівнянь
і складається з наступного:

\begin{itemize}
\tightlist
\item
  \begin{enumerate}
  \def\labelenumi{\arabic{enumi})}
  \tightlist
  \item
    знайти загальний розв'язок \texttt{y\_n} однорідного рівняння;
  \end{enumerate}
\item
  \begin{enumerate}
  \def\labelenumi{\arabic{enumi})}
  \setcounter{enumi}{1}
  \tightlist
  \item
    знайти часткове розв'язок \texttt{y\_p} гетерогенного рівняння;
  \end{enumerate}
\item
  \begin{enumerate}
  \def\labelenumi{\arabic{enumi})}
  \setcounter{enumi}{2}
  \tightlist
  \item
    Підставити \texttt{y\_h\ +\ y\_p} у початкову умову та визначити
    константи.
  \end{enumerate}
\end{itemize}

\texttt{Ми\ використаємо\ цей\ метод\ для\ розв\textquotesingle{}язання\ нашої\ проблеми.\ У\ нашій\ задачі\ \textquotesingle{}P\_1\textquotesingle{}\ загальний\ розв\textquotesingle{}язок\ однорідного\ рівняння\ \textquotesingle{}\textbackslash{}Delta\ u=0\textquotesingle{}\ має\ вигляд\ із\ розділеними\ змінними}

\begin{verbatim}
\begin{aligned} u_{tt} &= c^2 u_{xx}, & 0 < x < 1, & 0 < t < \infty, \\ u(0, t) &= g_1(t), & 0 < t < \infty, \\ u(1, t) &= g_2(t), & 0 < t < \infty, \\ (HY) &\begin{cases} u(x, 0) &= f(x), \\ u_t(x, 0) &= g(x), \end{cases} & 0 \leqslant x \leqslant 1, \end{aligned}
\end{verbatim}

Часткове (лише одне) розв'язок неоднорідного рівняння

Ми розглядаємо у формі

\begin{verbatim}
u_x(0, t) = 0, \quad u_x(L, t) = 0.
\end{verbatim}

\begin{verbatim}
Правильні частини форми r^n , r^n \cos{(n\theta)} , r^n \sin{(n\theta)} ведуть до рішень Ar^{n+2} , Br^{n+2} \cos{(n\theta)} Cr^{n+2} \sin{(n\theta)} відповідно.
\end{verbatim}

\begin{verbatim}
Підставляючи функції u_p(r, \theta) у гетерогенне рівняння, отримуємо A = -1/32, B = -1/24.
\end{verbatim}

Отже,

\begin{verbatim}
u_{x}(1, t) = \frac{1}{k} v(t) - (k - \text{модуль Юнга}).
\end{verbatim}

\begin{verbatim}
Щоб розв'язати задачу, P_1 має зробити останній крок — підставити загальне розв'язання u(r, \theta) = u_h(r, \theta) + u_p(r, \theta) на граничні умови u(1, \theta) = 0 і знайти коефіцієнти a_n і b_n .
\end{verbatim}

Після цього ми отримуємо

\begin{verbatim}
u_{x}(0, t) = \frac{h}{T} u(0, t),
\end{verbatim}

Отже, \texttt{a\_0=1/32\ ,\ a\_2=1/24},і всі інші \texttt{a\_n} і
\texttt{b\_n} дорівнюють нулю. Отже, рішення проблеми \texttt{P\_1}
знайдено:

\begin{verbatim}
u_{x}(L, t) = -\frac{h}{T} u(L, t).
\end{verbatim}

\begin{verbatim}
Функція u_1(r, \theta) називається пертурбацією першого порядку до u_4(r, \theta) . Додавши u_0 і u_1 , отримуємо новий підхід до розв'язання задачі (46.2):
\end{verbatim}

(46.5)

\begin{verbatim}
u_{x}(0, t) - \frac{h}{T}u(0, t) = 0,
\end{verbatim}

\begin{verbatim}
Щоб знайти збурення наступного порядку u_{z}(r, \theta) , потрібно розв'язати задачу
\end{verbatim}

\begin{verbatim}
u_{x}(L, t) + \frac{h}{T}u(L, t) = 0.
\begin{cases} \Delta u_{1} = -2u_{0}u_{1}, \\ u_{2}(1, \theta) = 0, \end{cases}
\end{verbatim}

\begin{verbatim}
де u_{\phi}(r,\theta) і u_{1}(r,0) вже відомі, раніше знайдені функції. Звісно, виконувати всі ці алгебраїчні перетворення без допомоги комп'ютера надзвичайно складно. На щастя, вираз (46.5) є майже точним рішенням нашої проблеми. Насправді, якщо підставити цей вираз зліва від \Delta u + u^{2} нашого нелінійного рівняння, отримаємо вираз, який майже всюди відрізняється від нуля в межах кола 0 < r < 1.
\end{verbatim}

Окрім розв'язання нелінійних задач, теорія збурень може
використовуватися для розв'язання задач у областях з неправильними
межами (звісно, якщо нерегулярність не дуже велика). Давайте розглянемо
простий приклад.

\section{Приклад задачі з порушеною
межею}\label{ux43fux440ux438ux43aux43bux430ux434-ux437ux430ux434ux430ux447ux456-ux437-ux43fux43eux440ux443ux448ux435ux43dux43eux44e-ux43cux435ux436ux435ux44e}

\begin{verbatim}
Порушення можна порушити не лише рівняння з похідними похідними. Можливо знайти розв'язок рівняння Лапласа всередині деформованого кола, збуривши розв'язок рівняння Лапласа всередині кругової області. Наприклад, припустимо, ми хочемо знайти потенціал у площі, обмеженій кривою r=1+\frac{1}{4}\sin\theta (у полярних координатах). Нехай потенціал u на межі заданий. Тобто, нам потрібно вирішити проблему
\end{verbatim}

(46.6)

\begin{verbatim}
u \left( 1 + \frac{1}{4} \sin \theta, \ \theta \right) = \cos \theta, \quad 0 \le \theta \le 2\pi.
\begin{aligned} u_x(0, T) &= \frac{h}{T} [u(0, t) - \theta_1(t)], \\ u_x(L, T) &= -\frac{h}{T} [u(L, t) - \theta_2(t)]. \end{aligned}
u\left(1+\frac{1}{4}\sin\theta,\ \theta\right)=\cos\theta
u(1+\epsilon\sin\theta, \theta) = \cos\theta, \quad 0 \le \epsilon \le 1/4.
\end{verbatim}

Этим мы завершаем обсуждение наиболее важных граничных условий,
связанных с гиперболическими задачами. Следующие несколько лекций мы
посвятим решению задач, содержащих полобные граничные условия.

\section{ЗАМЕЧАНИЯ}\label{ux437ux430ux43cux435ux447ux430ux43dux438ux44f-4}

\begin{enumerate}
\def\labelenumi{\arabic{enumi}.}
\tightlist
\item
  В лекции не рассматривалось граничное условие еще одного типа, когда
  струна подвергается действию силы, пропорциональной скорости и
  направленной в противоположном направ-
\end{enumerate}

лении. Граничное условие такого типа (на левом конце) записывается в
виде

\begin{verbatim}
f(x+h) = f(x) + f'(x)h + f''(x)\frac{h^2}{2!} + \dots,
\end{verbatim}

\begin{enumerate}
\def\labelenumi{\arabic{enumi}.}
\setcounter{enumi}{1}
\tightlist
\item
  Нелинейное упругое закрепление левого конца приводит к граничному
  условию вида
\end{enumerate}

\begin{verbatim}
u(1+\varepsilon\sin\theta,\ \theta)=u(1,\ \theta)+u_r(1,\ \theta)(\varepsilon\sin\theta)+u_{rr}(1,\ \theta)\frac{(\varepsilon\sin\theta)^2}{21}+\ldots
\end{verbatim}

где \texttt{\textbackslash{}varphi(u)} --- произвольная функция величины
u. Например, граничное условие

\begin{verbatim}
\Delta u = 0, \quad 0 < r < 1 + \varepsilon \sin \theta,
\end{verbatim}

говорит о том, что возвращающая сила на левом конце пропорциональна кубу
смещения (а не первой степени и, как это будет согласно закону Гука в
линейном случае).

\begin{enumerate}
\def\labelenumi{\arabic{enumi}.}
\setcounter{enumi}{2}
\tightlist
\item
  Если к нижнему концу пружины, совершающей продольные колебания,
  прикрепить массу \emph{m}, то граничное условие следует задавать в
  виде
\end{enumerate}

\begin{verbatim}
Коли всі розрахунки будуть виконані, ми внесемо \varepsilon=1/4 у кінцеві результати. Звісно, здається, що цю задачу не можна легко розв'язати, але її можна звести до послідовності задач, у кожній з яких визначена одна з функцій u_{\mathrm{ }e} , u_{\mathrm{ }e} , ..., що включені у розклад
\end{verbatim}

\section{ЗАДАЧИ}\label{ux437ux430ux434ux430ux447ux438-2}

Опираясь на интуитивные представления о граничных условиях различных
типов, изобразить в общих чертах решение задачи

Після підстави ряду в (467) отримуємо таку послідовність задач для
визначення функцій \texttt{u\_a} , \texttt{u\_1} , \ldots;

\begin{verbatim}
\begin{cases} \Delta u_{0} = 0, & 0 < r < 1 \text{ (внутри круга),} \\ u_{0}(1, \theta) = \cos \theta, & u_{0}(r, \theta) = r \cos \theta, \end{cases}
\end{verbatim}

\begin{verbatim}
\begin{array}{ll} (\mathrm{Y} \, \Pi) & u_{tt} = c^2 u_{xx}, \quad 0 < x < 1, \quad 0 < t < \infty, \\ (\mathrm{\Gamma} \, \mathrm{Y}) & \left\{ \begin{array}{ll} u \, (0, \, t) = 0, \\ u \, (1, \, t) = \sin t, \end{array} \end{array}
\end{verbatim}

\begin{verbatim}
\begin{cases} \Delta u_{1} = 0, & 0 < r < 1 \text{ (внутри круга),} \\ u_{1}(1, \theta) = -\sin \theta \frac{\partial u_{0}(1, \theta)}{\partial r} = -\sin \theta \cos \theta \end{cases}
u = u_0 + \frac{1}{4}u_1 + \left(\frac{1}{4}\right)^2 u_1 + \cdots
u_u + \frac{1}{4} u_t
\begin{array}{ll} u \, (x, \, 0) = 0, \\ u_{t} \, (x, \, 0) = 0, \end{array}
u_t = (1+x) u_{xx}, \\ u(x, 0) = q(x), \quad -\infty < x < \infty,
\end{verbatim}

\begin{verbatim}
\right. & 0 \leqslant x \leqslant 1, \end{array}
\end{verbatim}

\begin{verbatim}
u_1 = (1 + \varepsilon_1) u_{xx}
u = u_1 + \varepsilon u_1 + \varepsilon^2 u_2 + \cdots
P_{0}
\end{verbatim}

в различные моменты времени. Могли бы вы угадать решение этой задачи?

\begin{enumerate}
\def\labelenumi{\arabic{enumi}.}
\setcounter{enumi}{2}
\tightlist
\item
  Во что перейдет граничное условие
\end{enumerate}

Отметим, что в первых двух задачах коэффициенты постоянны. Читатель
должен понимать, что параметр в должен быть достаточно мал, ниаче
бескопечный ряд может оказаться расходящимся

\section{ЗАДАЧИ}\label{ux437ux430ux434ux430ux447ux438-3}

\begin{itemize}
\tightlist
\item
  \begin{enumerate}
  \def\labelenumi{\arabic{enumi}.}
  \tightlist
  \item
    Подставьте разложение (46.8) в задачу (46.7) и получите
    последовательность задач
    \texttt{P\_0,\textbackslash{}\ P\_1,\textbackslash{}\ P\_2,\textbackslash{}\ \textbackslash{}dots}
    . 2. Покажите, что нелинейную задачу
  \end{enumerate}
\end{itemize}

npit

\begin{verbatim}
- a)   h \rightarrow \infty  ,
- 6)   h \rightarrow 0  .
\end{verbatim}

Согласуется ли результат с вашей интунцией?

\begin{enumerate}
\def\labelenumi{\arabic{enumi}.}
\setcounter{enumi}{3}
\tightlist
\item
  Изобразить в общих чертах решение задачи
\end{enumerate}

(УЧП)

можна звести до послідовності лінійних задач \texttt{P\_\{o\}} ,
\texttt{P\_\{i\}} , \texttt{P\_\{z\}} .

\begin{itemize}
\tightlist
\item
  \begin{enumerate}
  \def\labelenumi{\arabic{enumi}.}
  \setcounter{enumi}{2}
  \tightlist
  \item
    Замініть нелінійну задачу 2 на приблизне розв'язання (46,5) і
    оцініть її точність.
  \end{enumerate}
\item
  \begin{enumerate}
  \def\labelenumi{\arabic{enumi}.}
  \setcounter{enumi}{3}
  \tightlist
  \item
    Розв'яжіть задачу \texttt{P\_1} з прикладу з збуреною межою і
    перевірте, наскільки добре є наближене розв'язання
  \end{enumerate}
\end{itemize}

(ГУ)

задовольняє відношення

\texttt{\$\$\ \textquotesingle{}u\textbackslash{}left(1\ +\ \textbackslash{}frac\{1\}\{4\}\textbackslash{}sin\textbackslash{}theta,\ \textbackslash{}theta\textbackslash{}right)\ =\ \textbackslash{}cos\textbackslash{}theta.\textquotesingle{}\ \#\ Лекція\ 47}

\end{document}
